\documentclass[../main]{subfiles}
\begin{document}

\chapter{Segundo Principio de la Termodinámica}

\section{Introducción al Segundo Principio de la Termodinámica}

El Segundo Principio de la Termodinámica es una de las leyes fundamentales que rigen el comportamiento de la energía en sistemas termodinámicos. Esta ley, formulada en diversas maneras, establece restricciones sobre la dirección del flujo de calor y la transformación de la energía térmica en trabajo mecánico.

\subsection{Enunciado de Kelvin-Planck}

Una de las formulaciones del Segundo Principio es el enunciado de Kelvin-Planck, que establece que es imposible construir una máquina térmica que opere en un ciclo y tenga como único efecto la extracción de una cantidad de calor de un cuerpo y la completa conversión de ese calor en trabajo. En otras palabras, no es posible construir una máquina que tome calor de una sola fuente y lo convierta completamente en trabajo mecánico sin dejar algún efecto secundario.

\subsection{Enunciado de Clausius}

El enunciado de Clausius del Segundo Principio se basa en el concepto de la entropía. Establece que es imposible que una máquina cíclica opere sin pérdida y produzca trabajo útil, tomando calor de un solo cuerpo y transmitiéndolo a otro a temperatura más alta.

\subsection{Entropía y el Segundo Principio}

La entropía es una medida del desorden o la aleatoriedad de un sistema. El Segundo Principio se puede entender en términos de la entropía, ya que afirma que en un sistema aislado, la entropía tiende a aumentar con el tiempo.

\info{Segundo Principio de la Termodinámica}{
El Segundo Principio de la Termodinámica establece restricciones fundamentales sobre la dirección del flujo de calor y la transformación de la energía térmica en trabajo mecánico. 

Una formulación importante es el enunciado de Kelvin-Planck, que establece que es imposible construir una máquina térmica que opere en un ciclo y tenga como único efecto la extracción de una cantidad de calor de un cuerpo y la completa conversión de ese calor en trabajo. 

Otra formulación es el enunciado de Clausius, que se basa en el concepto de entropía y establece que es imposible que una máquina cíclica opere sin pérdida y produzca trabajo útil, tomando calor de un solo cuerpo y transmitiéndolo a otro a temperatura más alta.
}

% Aquí se pueden agregar más secciones explicativas y fórmulas relacionadas al Segundo Principio

\actividad{Cuestionario sobre el Segundo Principio de la Termodinámica}

\begin{enumerate}
    \item ¿Cuáles son las formulaciones principales del Segundo Principio de la Termodinámica?
    \item ¿Qué establece el enunciado de Kelvin-Planck?
    \item Explique el concepto de entropía y su relación con el Segundo Principio.
    % Preguntas adicionales
\end{enumerate}

\actividad{Problemas relacionados con el Segundo Principio de la Termodinámica}

\begin{enumerate}
    \item Una máquina térmica opera entre dos fuentes de calor a temperaturas $T_1 = 500K$ y $T_2 = 300K$. Si la eficiencia de la máquina es del $40\%$, ¿cuánto trabajo realiza la máquina si absorbe $5000J$ de calor de la fuente caliente?
    \item Un refrigerador tiene una eficiencia de $25\%$. Si necesita extraer $4000J$ de calor de un espacio a $300K$, ¿cuánto trabajo debe realizar para lograrlo?
    % Otros problemas
\end{enumerate}

\end{document}
